\chapter{Project: A Transformer Language Model from Scratch (CS336 Assignment 1)}\label{ch:proj1}

Part~VII turns the book into a set of projects: Stanford's CS336, \emph{Language Modeling
from Scratch} (Spring 2026), publishes five implementation-heavy assignments, and each of the
following five chapters is a guide to one of them. The guides tell you how to set up, which
chapters and labs teach the ideas each problem needs, where the traps are, how to design and
log the experiments, and how to scale the work down when you do not have the course's GPUs.
They deliberately stop short of solutions.

\begin{keyidea}[How to use the project guides (read this first)]
The CS336 assignments are designed to be done \emph{by you}. The course's policy, which the
assignment repositories also include as \code{AGENTS.md}/\code{CLAUDE.md} instructions for AI
tools, allows help with concepts, documentation and debugging strategy but not with
implementing assignment components. These guides follow the same line:
\begin{itemize}
  \item every assignment is described in our own words, and the official handout is the
        specification. When the two disagree, the handout wins;
  \item there is no solution code for any assignment problem and no final answers to the
        written or accounting questions. Where a problem needs a method, we demonstrate it on a
        \emph{different} example;
  \item the book's own chapters and labs implement many of the same \emph{concepts} (a BPE
        trainer, an AdamW, a Transformer block) against the book's \code{llmsys} API. If you are
        taking the course for credit, do not copy them into your submission: write your own,
        and use the book's versions only to understand what correct behaviour looks like.
\end{itemize}
If you are working through the assignments on your own, the same advice applies for a different
reason: the value of these assignments is in the debugging you do along the way.
\end{keyidea}

% =====================================================================
\section{What you build}\label{sec:proj1:overview}\covers{CS336 A1 \S1; CS336 L1}

Assignment~1 asks you to build, with almost no library help, everything needed to train a
modern decoder-only language model, and then to train and study such models. In our words:

\begin{enumerate}
  \item a \textbf{byte-level BPE tokenizer}: a trainer that learns merges from a text file and a
        tokenizer that encodes and decodes with them, including special tokens and streaming input;
  \item a \textbf{pre-norm Transformer LM} with RMSNorm, SwiGLU feed-forward layers, rotary
        position embeddings and causal multi-head attention, built from your own \code{Linear}
        and \code{Embedding} modules;
  \item the \textbf{training machinery}: cross-entropy, an AdamW optimizer, a cosine schedule with
        warmup, gradient clipping, a memory-mapped data loader, checkpointing and a training script;
  \item \textbf{decoding} with temperature and nucleus sampling;
  \item \textbf{experiments} on TinyStories and an OpenWebText sample: learning-rate and
        batch-size studies, four architecture ablations, a time-limited leaderboard run, and a
        written report with an experiment log.
\end{enumerate}

The ``from scratch'' rule is concrete: from \code{torch.nn}, \code{torch.nn.functional} and
\code{torch.optim} you may use only \code{nn.Parameter}, the container classes
(\code{Module}, \code{ModuleList}, \code{Sequential}, \ldots) and the \code{Optimizer} base class.
Tensor operations, \code{einsum}, \code{torch.nn.init.trunc\_normal\_} and similar primitives are
allowed. Check the handout for the exact list.

\begin{figure}[t]
\centering
\resizebox{\linewidth}{!}{%
\begin{tikzpicture}[
  comp/.style={draw, rounded corners=3pt, minimum height=0.9cm, align=center, font=\small, line width=0.6pt, text width=2.6cm},
  ref/.style={font=\scriptsize, text=srcgray, align=center},
  arr/.style={-{Stealth[length=2mm]}, line width=0.6pt}]
\node[comp, fill=fdata!18, draw=keygreen] (tok) at (0,0) {BPE trainer +\\tokenizer};
\node[ref, below=1pt of tok] {Ch.~3, Lab~3; \S\ref{sec:proj1:bpe}};
\node[comp, fill=fdata!18, draw=keygreen] (ids) at (0,-2.2) {uint16 token\\files};
\node[ref, below=1pt of ids] {Ch.~3, Ch.~9};
\node[comp, fill=fcompute!18, draw=sysblue] (mods) at (4.2,0) {Linear, Embedding,\\RMSNorm, SwiGLU};
\node[ref, below=1pt of mods] {Ch.~6, Ch.~7};
\node[comp, fill=fcompute!18, draw=sysblue] (attn) at (4.2,-2.2) {RoPE, softmax,\\causal MHA};
\node[ref, below=1pt of attn] {Ch.~4, Ch.~6};
\node[comp, fill=fcompute!18, draw=sysblue] (lm) at (8.4,-1.1) {pre-norm block\\$\to$ Transformer LM};
\node[ref, below=1pt of lm] {Ch.~5, Ch.~6};
\node[comp, fill=fmemory!20, draw=napkin] (opt) at (12.6,0) {cross-entropy, AdamW,\\schedule, clipping};
\node[ref, below=1pt of opt] {Ch.~9, Ch.~10};
\node[comp, fill=fmemory!20, draw=napkin] (loop) at (12.6,-2.2) {loader, checkpoints,\\training script};
\node[ref, below=1pt of loop] {Ch.~9, Lab~9};
\node[comp, fill=fcomm!16, draw=labpurple] (exp) at (16.8,-1.1) {experiments,\\ablations,\\leaderboard};
\node[ref, below=1pt of exp] {Ch.~7, Ch.~9, \S\ref{sec:proj1:experiments}};
\draw[arr] (tok)--(ids); \draw[arr] (mods)--(lm); \draw[arr] (attn)--(lm);
\draw[arr] (lm)--(opt); \draw[arr] (lm)--(loop); \draw[arr] (ids.east) to[out=-15,in=200] (loop.west);
\draw[arr] (opt)--(exp); \draw[arr] (loop)--(exp);
\end{tikzpicture}}
\caption{The components of Assignment~1, in dependency order, with the chapters that teach the
underlying ideas. Each box is several of the handout's problems; \cref{tab:proj1:map} gives the
problem-by-problem map.}
\label{fig:proj1:components}
\end{figure}

% =====================================================================
\section{Setting up}\label{sec:proj1:setup}\covers{CS336 A1 \S1}

\paragraph{Repository and environment.} The assignment lives in a public GitHub repository
(\url{https://github.com/stanford-cs336/assignment1-basics}) together with the handout PDF. It
uses \texttt{uv} to pin Python (3.12 or 3.13) and the dependencies: a specific PyTorch minor
version, \texttt{einops}, \texttt{einx}, \texttt{jaxtyping}, \texttt{regex}, \texttt{tiktoken}
(used by the tests as a reference tokenizer), \texttt{wandb}, \texttt{pytest} and a few tools.

\begin{lstlisting}[style=shell]
git clone https://github.com/stanford-cs336/assignment1-basics && cd assignment1-basics
pip install uv            # or the standalone installer
uv run pytest             # first run creates .venv; every test fails with NotImplementedError
uv run python -c "import torch; print(torch.__version__, torch.cuda.is_available())"
python <book>/code/projects/a1/check_env.py     # machine facts + tips (Chapter 32 tools)
\end{lstlisting}

\paragraph{How your code meets the tests.} You write all your code in the package directory
(\texttt{cs336\_basics/}), structured however you like. The tests never import it directly.
They call functions in \texttt{tests/adapters.py}, each of which currently raises
\code{NotImplementedError}. You fill each adapter with a few lines of \emph{glue}: construct your
module, load the weights the test passes in (usually via \code{load\_state\_dict}), and call it
(\cref{fig:proj1:adapters}). The rule is that adapters contain no substantive logic. That keeps
your package usable outside the tests, and keeps the tests stable across different internal
designs.

\begin{figure}[t]
\centering
\begin{tikzpicture}[
  b/.style={draw, rounded corners=2pt, minimum height=0.85cm, align=center, font=\small, text width=3.4cm, line width=0.6pt},
  arr/.style={-{Stealth[length=2mm]}, line width=0.6pt}]
\node[b, fill=black!5] (t) at (0,0) {\texttt{tests/test\_model.py}\\(do not edit)};
\node[b, fill=fmemory!20, draw=napkin] (a) at (5,0) {\texttt{tests/adapters.py}\\glue only};
\node[b, fill=fcompute!18, draw=sysblue] (p) at (10,0) {\texttt{cs336\_basics/}\\your modules};
\node[b, fill=black!5] (s) at (0,-1.9) {\texttt{tests/\_snapshots/*.npz}\\\texttt{tests/fixtures/}};
\draw[arr] (t) -- node[above,font=\scriptsize]{calls \code{run\_*}} (a);
\draw[arr] (a) -- node[above,font=\scriptsize]{builds, loads, calls} (p);
\draw[arr, dashed] (s) -- node[left,font=\scriptsize]{reference outputs} (t);
\draw[arr, dashed] (p.south) |- node[pos=0.75,below,font=\scriptsize]{outputs compared to snapshots / tiktoken} (s.east);
\end{tikzpicture}
\caption{The adapter pattern used throughout CS336. The tests reach your code only through thin
adapter functions. Numerical tests compare against stored reference outputs (snapshots), and the
tokenizer tests compare against \texttt{tiktoken}'s GPT-2 tokenizer.}
\label{fig:proj1:adapters}
\end{figure}

A few properties of the test suite matter for how you work:

\begin{itemize}
  \item \textbf{Run subsets.} \code{uv run pytest -k test\_rmsnorm} runs one test. Work
        component by component, in the order of \cref{tab:proj1:map}.
  \item \textbf{Snapshot tests check numbers, not just shapes.} Model components are tested
        against stored outputs computed with specific weights the test loads into your module.
        Your parameter names and shapes must accept those weights (read the adapter docstrings for
        the expected shapes and orientation, e.g.\ $W$ stored as $d_\text{out}\times d_\text{in}$),
        and your computation must follow the handout's exact conventions. Two conventions where
        the book differs: RoPE pairs (rotate-half vs adjacent pairs, \cref{sec:modern:rope}) and
        mask polarity (\cref{sec:proj1:conventions}).
  \item \textbf{Some tests have resource limits.} One BPE test fails if training on a small corpus
        takes longer than a fixed number of seconds, and one streaming-tokenizer test runs under
        an address-space limit of about a megabyte above the current usage (enforced with
        \code{resource.setrlimit}, which works on Linux but not reliably on macOS). The submission
        script also runs every test with a per-test timeout.
  \item \textbf{The tokenizer tests use GPT-2's real vocabulary.} The fixtures store GPT-2's
        vocabulary and merges in the original serialised form, where bytes are mapped to printable
        Unicode characters. Converting between that form and raw \code{bytes} is part of your setup
        work.
\end{itemize}

\paragraph{Data.} The README lists download commands for TinyStories (train and validation
text files, documents separated by \texttt{<|endoftext|>}) and a sample of OpenWebText
(gzipped text). Both are single large plain-text files; budget tens of gigabytes of disk for
the text and the \code{uint16} token arrays you will create from them, and keep them out of your
submission zip (the submission script has an exclusion list).

\paragraph{Compute.} The handout states budgets in B200-hours, and the full TinyStories run is
designed to take well under an hour on a B200. It also gives a low-resource path: a reduced
token budget that trains a fluent-ish TinyStories model on an Apple-silicon laptop (MPS) or a CPU
in about an hour, with a correspondingly relaxed loss target. \Cref{sec:proj1:outside} discusses
hardware substitutions.

% =====================================================================
\section{Problem-by-problem map}\label{sec:proj1:map}

\Cref{tab:proj1:map} lists every problem in the handout, what it asks for in a phrase, and where
this book develops the needed ideas. ``W'' marks written answers, ``C'' code, ``E'' experiments.

{\small
\begin{longtable}{@{}p{0.27\linewidth}p{0.06\linewidth}p{0.30\linewidth}p{0.29\linewidth}@{}}
\caption{Assignment~1 problems mapped to the book.}\label{tab:proj1:map}\\
\toprule
\textbf{Problem} & & \textbf{What it asks (paraphrased)} & \textbf{Where the book helps}\\
\midrule\endfirsthead
\toprule
\textbf{Problem} & & \textbf{What it asks} & \textbf{Where the book helps}\\
\midrule\endhead
\multicolumn{4}{@{}l}{\emph{Tokenization (handout \S2)}}\\
\code{unicode1}, \code{unicode2} & W & reason about code points, the NUL character, UTF-8 vs UTF-16/32, invalid byte sequences & \cref{sec:text:wild}; \cref{sec:proj1:unicode}\\
\code{train\_bpe} & C & byte-level BPE trainer with special tokens as hard boundaries, deterministic tie-breaking, fast enough for a timed test & \cref{sec:text:bpe}, Lab~3; \cref{sec:proj1:bpe}\\
\code{train\_bpe\_tinystories}, \code{train\_bpe\_expts\_owt} & E/W & train 10K and 32K vocabularies; time, memory, profile, longest token & \cref{sec:proj1:bpe} (parallelism, profiling); \cref{sec:text:pipeline}\\
\code{tokenizer} & C & encoder/decoder class with special tokens, file loading and lazy \code{encode\_iterable} & \cref{sec:text:bpe,sec:text:special}; \cref{sec:proj1:stream}\\
\code{tokenizer\_experiments} & E/W & compression ratios, cross-domain behaviour, throughput, why \code{uint16} & \cref{sec:text:pipeline}, \cref{sec:optim:data}; \cref{sec:proj1:stream}\\
\multicolumn{4}{@{}l}{\emph{Architecture (handout \S3)}}\\
\code{linear}, \code{embedding} & C & bias-free linear and embedding modules with truncated-normal init & \cref{sec:archdesign:einsum,sec:archdesign:init}; \cref{sec:text:tokemb}\\
\code{rmsnorm} & C & RMSNorm with fp32 upcast & \cref{sec:modern:rmsnorm}\\
\code{positionwise\_feedforward} & C & SwiGLU FFN, $\dff\approx\tfrac83 d$ rounded to a multiple of 64 & \cref{sec:modern:swiglu}, \cref{sec:archdesign:mlp}\\
\code{rope} & C & RoPE module with precomputed cos/sin buffers and explicit token positions & \cref{sec:modern:rope}, \cref{sec:archdesign:pos}; \cref{sec:proj1:conventions}\\
\code{softmax}, \code{scaled\_dot\_product\_attention} & C & stable softmax; attention with optional boolean mask and arbitrary batch dims & \cref{sec:attention:simple,sec:attention:causal}, \cref{sec:optim:loss}\\
\code{multihead\_self\_attention} & C & causal MHA with RoPE on $Q,K$, three projections & \cref{sec:attention:mha}, \cref{sec:modern:gqa}\\
\code{transformer\_block}, \code{transformer\_lm} & C & pre-norm block and full LM & \cref{sec:gpt:block}, \cref{sec:modern:block,sec:modern:model}\\
\code{transformer\_accounting} & W & parameter count, matmul FLOPs and their breakdown for several model sizes and context lengths & \cref{sec:gpt:accounting}, Lab~1, \cref{sec:optim:account}\\
\multicolumn{4}{@{}l}{\emph{Training (handout \S4--5)}}\\
\code{cross\_entropy} & C & stable cross-entropy averaged over batch dimensions & \cref{sec:pretrain:loss}, \cref{sec:optim:loss}\\
\code{learning\_rate\_tuning} & E/W & toy SGD at several learning rates & \cref{sec:optim:api} (pitfall box)\\
\code{adamw} & C & AdamW as an \code{Optimizer} subclass & \cref{sec:optim:api,sec:optim:adam}, Lab~9\\
\code{adamw\_accounting} & W & peak memory decomposition, max batch size, FLOPs of a step, training time at given MFU & \cref{sec:optim:account,sec:optim:cost}, \cref{sec:pretrain:memory}\\
\code{learning\_rate\_schedule}, \code{gradient\_clipping} & C & warmup + cosine; global-norm clipping & \cref{sec:optim:sched,sec:optim:clip}\\
\code{data\_loading}, \code{checkpointing} & C & random batches from a token array; save/load model + optimizer + iteration & \cref{sec:optim:data,sec:optim:ckpt}\\
\code{training\_together} & C & configurable training script with memmap data, checkpoints, logging & \cref{sec:optim:loop,sec:optim:explog}, \cref{sec:pretrain:trainfn}\\
\code{decoding} & C & sampling with temperature and top-$p$ until EOT or a length limit & \cref{sec:pretrain:decoding}, \cref{sec:scaling:topp}\\
\multicolumn{4}{@{}l}{\emph{Experiments (handout \S7)}}\\
\code{experiment\_log} & C/W & tracking infrastructure (step and wall clock) and a log of what you tried & \cref{sec:optim:explog}; \cref{sec:proj1:experiments}\\
\code{learning\_rate}, \code{batch\_size\_experiment} & E & LR sweep incl.\ a divergent run; batch-size sweep with re-tuned LRs & \cref{sec:optim:debug}, \cref{ch:scalinglaws}; \cref{sec:proj1:experiments}\\
\code{generate} & E/W & sample text, comment on fluency & \cref{sec:pretrain:decoding}\\
\code{layer\_norm\_ablation}, \code{pre\_norm\_ablation}, \code{no\_pos\_emb}, \code{swiglu\_ablation} & E & remove norms; post-norm; NoPE; SiLU vs SwiGLU at matched parameters & \cref{sec:archdesign:norm,sec:archdesign:pos,sec:archdesign:mlp,sec:archdesign:ablate}, Lab~7\\
\code{main\_experiment} & E/W & same recipe on OpenWebText; interpret the loss gap & \cref{sec:proj1:owt}, \cref{ch:data}\\
\code{leaderboard} & E & best validation loss within a fixed wall-clock budget & \cref{sec:proj1:leaderboard}\\
\bottomrule
\end{longtable}}

% =====================================================================
\section{Background the assignment needs that earlier chapters only touch on}\label{sec:proj1:background}

Most of the assignment rests on Chapters 3--10. This section fills the gaps the handout relies
on most: Unicode details, the exact pre-tokenisation regex, scaling BPE training to a large file, and
streaming encoders. It also lists the conventions where the book and the handout differ.

\subsection{Unicode, code points and UTF-8}\label{sec:proj1:unicode}\covers{CS336 A1 \S2.1--2.2}

Unicode assigns each character an integer \emph{code point} (\code{ord}/\code{chr} in Python);
the standard defines roughly 160,000 characters across a space of $17\times2^{16}$ code points.
Some code points are control characters that print as nothing visible. An \emph{encoding}
turns code points into bytes. UTF-8 is variable-length:

\begin{center}\small
\begin{tabular}{@{}lll@{}}
\toprule
code point range & bytes & bit pattern\\
\midrule
U+0000--U+007F & 1 & \texttt{0xxxxxxx} (ASCII unchanged)\\
U+0080--U+07FF & 2 & \texttt{110xxxxx 10xxxxxx}\\
U+0800--U+FFFF & 3 & \texttt{1110xxxx 10xxxxxx 10xxxxxx}\\
U+10000--U+10FFFF & 4 & \texttt{11110xxx 10xxxxxx 10xxxxxx 10xxxxxx}\\
\bottomrule
\end{tabular}
\end{center}

Three consequences matter for tokenizers. First, a byte-level vocabulary needs only 256 base
symbols and can never meet an out-of-vocabulary input. Second, a single byte is not a character:
continuation bytes (\texttt{10xxxxxx}) are meaningless on their own, so any code that decodes
token-by-token or byte-by-byte must buffer incomplete sequences (\cref{sec:text:wild} builds an
incremental detokenizer for exactly this reason). Third, not every byte sequence is valid UTF-8:
a lead byte must be followed by the right number of continuation bytes, and some lead bytes never
occur at all. A decoder therefore needs an error policy. The handout asks for replacement with
U+FFFD, which Python provides through \code{bytes.decode("utf-8", errors="replace")}. Explore
the alternatives yourself by encoding the same strings in UTF-16 and UTF-32 and comparing lengths
and byte patterns: that is the evidence the written questions ask for.

\subsection{Pre-tokenisation with the GPT-2 regex}\covers{CS336 A1 \S2.4}

The book's \code{ByteBPETokenizer} (\cref{sec:text:bpe}) uses a simplified ASCII-only pattern.
The assignment uses the GPT-2 pattern, which needs Unicode character classes:
\begin{lstlisting}[style=plain]
'(?:[sdmt]|ll|ve|re)| ?\p{L}+| ?\p{N}+| ?[^\s\p{L}\p{N}]+|\s+(?!\S)|\s+
\end{lstlisting}
Python's built-in \code{re} does not support \verb|\p{L}| (``any letter in any script'') or
\verb|\p{N}| (``any number''); the third-party \code{regex} module does. Read the alternatives
left to right: English contractions; an optional space followed by a run of letters; the same for
digits; the same for other non-space symbols; whitespace that is \emph{not} followed by a
non-space character (the negative lookahead \verb|\s+(?!\S)| leaves the last space of a run
attached to the next word, so ``\verb*|a   b|'' yields \verb*|a|, \verb*|  |, \verb*| b|); and
finally any remaining whitespace. Use \code{finditer} rather than \code{findall} when counting, so
you never hold a list of every pre-token of a multi-gigabyte file in memory.

\subsection{Scaling BPE training to a large file}\label{sec:proj1:bpe}\covers{CS336 A1 \S2.4--2.5}

\Cref{sec:text:bpe} derives BPE, and Lab~3 measures how an incremental trainer beats the naive
one. At the scale of TinyStories (gigabytes of text) and the OpenWebText sample, two system-level
issues dominate.

\paragraph{Pre-tokenisation is embarrassingly parallel, if you cut the file in the right
places.} Counting pre-tokens is a map--reduce: each worker counts pre-tokens in a slice of the
file, and the counters are summed. The only subtlety is the slice boundaries. A cut in the middle of
a word would create two fake pre-tokens. Special tokens already define hard boundaries (no merge
may cross a document separator), so cut \emph{at} occurrences of the separator. Pick
evenly spaced byte offsets, then move each one forward to the next separator
(\cref{fig:proj1:chunks}). The assignment repository provides a helper that computes such
boundaries, and you may use it as is. Then:
\begin{itemize}
  \item each worker opens the file itself and reads only its byte range, so you send offsets to
        the workers rather than text;
  \item within a chunk, split on \emph{all} special tokens (an alternation of
        \code{re.escape}d strings) before applying the regex, and do not count the special tokens;
  \item return a \code{Counter} of pre-token \emph{byte strings}; merging $k$ counters is cheap
        compared to the regex work.
\end{itemize}

\begin{figure}[t]
\centering
\begin{tikzpicture}[x=0.75cm, y=0.6cm]
\fill[fdata!12] (0,0) rectangle (20,1);
\foreach \x in {2.3,5.1,7.7,9.2,12.6,14.1,16.9,18.4} {\fill[warnred!50] (\x,0) rectangle +(0.25,1);}
\draw[black!50] (0,0) rectangle (20,1);
\node[font=\scriptsize, anchor=east] at (-0.1,0.5) {file};
\foreach \x/\l in {5/$b_1$,10/$b_2$,15/$b_3$} {\draw[dashed, black!60] (\x,-0.3)--(\x,1.3); \node[font=\scriptsize] at (\x,1.65) {\l};}
\foreach \x/\y in {5/5.1,10/12.6,15/16.9} {\draw[-{Stealth[length=1.6mm]}, sysblue, line width=0.7pt] (\x,-0.55) -- (\y,-0.55); \draw[sysblue, line width=0.9pt] (\y,-0.9)--(\y,0);}
\node[font=\scriptsize, sysblue, anchor=north] at (11.3,-0.9) {snap each guess forward to the next separator};
\draw[decorate, decoration={brace, amplitude=3pt, mirror}] (0,-1.8)--(5.1,-1.8) node[midway,below=3pt,font=\scriptsize]{worker 0};
\draw[decorate, decoration={brace, amplitude=3pt, mirror}] (5.1,-1.8)--(12.6,-1.8) node[midway,below=3pt,font=\scriptsize]{worker 1};
\draw[decorate, decoration={brace, amplitude=3pt, mirror}] (12.6,-1.8)--(16.9,-1.8) node[midway,below=3pt,font=\scriptsize]{worker 2};
\draw[decorate, decoration={brace, amplitude=3pt, mirror}] (16.9,-1.8)--(20,-1.8) node[midway,below=3pt,font=\scriptsize]{worker 3};
\node[font=\scriptsize, warnred, anchor=west] at (20.2,0.5) {\texttt{<|endoftext|>}};
\end{tikzpicture}
\caption{Parallel pre-tokenisation. Evenly spaced guesses $b_i$ are moved to the next document
separator, so no pre-token or merge can straddle two workers. Chunks have unequal sizes; use more
chunks than processes to balance the load.}
\label{fig:proj1:chunks}
\end{figure}

\paragraph{The merge loop is sequential, so make each iteration incremental.} After
pre-tokenisation the corpus is a dictionary from pre-token (a tuple of byte strings) to count,
typically with millions of distinct entries. Recounting all pairs after every merge costs
$O(\text{corpus})$ per merge. The standard fix is to maintain two indices: pair $\to$ count,
and pair $\to$ the set of pre-tokens that contain it. A merge then touches only the pre-tokens
that contain the merged pair, updating counts for their neighbouring pairs. Finding the maximum
can be a linear scan over the pair-count table (often fast enough) or a heap with lazy deletion.
Ties must be broken deterministically: the handout specifies preferring the lexicographically
greater pair, compared as \code{bytes}. Profile before optimising (\code{cProfile}, or
\code{py-spy} for a sampling profile of a running process): on large corpora the regex
pass usually dominates until you parallelise it, and the merge loop dominates afterwards.

\begin{napkinmath}[Is your trainer fast enough?]
Measure pre-tokenisation throughput on a 100\,MB slice with one process, in MB/s. A
single-threaded regex pass typically manages a few to a few tens of MB/s in pure Python. The
time for the full file is $\text{size}/(\text{throughput}\times P)$ with $P$ workers, plus the
merge loop, whose cost depends on the number of \emph{distinct} pre-tokens (not the file size)
times the number of merges times the average number of affected pre-tokens per merge. Extrapolate
both terms from small runs before launching a multi-hour job, and decide which one to optimise.
\end{napkinmath}

\subsection{Encoding: special tokens, streaming and serialisation}\label{sec:proj1:stream}\covers{CS336 A1 \S2.6--2.7}

\paragraph{Special tokens when encoding.} Split the input on special tokens first, emit their
ids directly, and BPE-encode only the text in between. When special tokens overlap (one is a
prefix or substring of another, such as a doubled end-of-text marker registered as its own
special token), build the alternation with the \emph{longest} tokens first so the regex prefers
them.

\paragraph{Streaming with \code{encode\_iterable}.} A generator that reads an iterable of
strings (a file handle yields lines) and yields ids lazily lets you tokenise a file that does not
fit in memory, in constant memory. The correctness condition is that the output must equal
\code{encode} on the concatenated text. That holds as long as you only cut the stream at points
where no pre-token or special token can straddle the cut. Lines are a convenient unit, but think
about what the regex does with the newline at the end of one line and the start of the next, and
about special tokens split across reads. Keep the per-chunk buffer bounded.

\paragraph{Throughput.} Real corpora repeat pre-tokens heavily, so caching pre-token $\to$ ids
(a plain dict or an LRU cache) typically speeds up encoding several-fold. Measure bytes per second
on a sample, extrapolate to your corpus (and to a 1\,TB corpus, to see why production tokenizers
are written in Rust), and use multiprocessing across documents for the full datasets.

\paragraph{Serialising vocabularies and merges.} Token strings are \code{bytes}, which are not
JSON-safe. Choose a lossless format: pickle, base64 strings in JSON, or GPT-2's printable
byte-to-Unicode mapping (the format of the GPT-2 fixtures in the test suite). Write a round-trip
test for whichever you choose. The arrays of token ids go to disk as \code{np.uint16}
(\cref{sec:optim:data}); assert that the vocabulary fits in that dtype, and memory-map the arrays
when training.

\subsection{Conventions that differ from this book}\label{sec:proj1:conventions}

The book's \code{llmsys} code and the assignment make different arbitrary choices in a few places.
The handout and the adapter docstrings are authoritative:

\begin{itemize}
  \item \textbf{RoPE pairing.} \code{llmsys.layers.apply\_rope} uses the ``rotate-half'' layout
        (dimension $i$ paired with $i+\dhead/2$), as Qwen3 and Llama checkpoints do. The handout
        defines rotation on \emph{adjacent} pairs $(2k-1,2k)$. Both are valid
        (\cref{sec:modern:rope} shows they are a fixed permutation of each other), but the snapshot
        test checks one specific layout. RoPE also takes explicit \emph{token positions} in the
        assignment's interface rather than an offset.
  \item \textbf{Mask polarity.} Our \code{causal\_mask} returns \code{True} where attention is
        \emph{forbidden}; the handout's boolean masks are \code{True} where attention is
        \emph{allowed}. Mixing the two silently inverts the causal mask, which shows up as a
        suspiciously low training loss.
  \item \textbf{Weight orientation and names.} Store $W\in\R^{d_\text{out}\times d_\text{in}}$ and
        apply it as $xW^\top$ (or with \code{einsum}, \cref{sec:archdesign:einsum}); the snapshot
        weights arrive in that orientation.
  \item \textbf{Initialisation.} The handout specifies truncated normals with given variances;
        \cref{sec:archdesign:init} explains where such rules come from.
  \item \textbf{AdamW.} Implement the handout's algorithm statement. Our \code{llmsys.optim.AdamW}
        follows PyTorch's formulation (\cref{lst:optim:adamw}), which is algebraically equivalent
        apart from where $\epsilon$ sits. For the test, follow the handout.
\end{itemize}

% =====================================================================
\section{Milestones and sanity checks}\label{sec:proj1:milestones}

Work in this order. Do not move on until each check passes.

\begin{enumerate}
  \item \textbf{Tokenizer.} The handout's toy example (a handful of words, whitespace
        pre-tokenisation) reproduces its stated merges when you train by hand. Round-trips are exact
        on Unicode, special tokens and empty strings. Ids match \code{tiktoken}'s GPT-2 encoding on the
        fixtures. \code{encode\_iterable(open(f))} equals \code{encode(open(f).read())}.
  \item \textbf{Components.} Each snapshot test passes individually. Before running the test, check
        shapes with toy inputs and an assertion or \code{jaxtyping} annotation per tensor.
  \item \textbf{Model at initialisation.} The loss on random data is close to $\ln V$
        (\cref{sec:optim:loss}): about 9.21 for a 10K vocabulary. The parameter count matches
        your hand count.
  \item \textbf{Optimizer.} Your AdamW and \code{torch.optim.AdamW} produce the same trajectory
        on a small problem (a useful private test, alongside the official one). Your schedule
        plotted over a full run looks like \cref{fig:optim:sched}.
  \item \textbf{Overfit a single batch} of TinyStories to near-zero loss in a few hundred steps.
  \item \textbf{Throughput.} Measure tokens/s and compare with $6N$ FLOPs per token and your
        device's peak (\cref{sec:optim:account}). If you are far below the handout's expected
        runtime, profile: the usual culprits are the data loader, evaluation on too many batches, a
        CPU sync every step, or unbatched attention.
  \item \textbf{Resume test.} Kill a run, resume from the checkpoint, and compare loss curves
        (Lab~9 automates this check for the book's loop).
\end{enumerate}

% =====================================================================
\section{Designing the experiments}\label{sec:proj1:experiments}\covers{CS336 A1 \S7.1--7.3}

\paragraph{Fix the token budget and derive steps from it.} The TinyStories experiments specify
the total number of tokens processed. For any batch size $B$ and context $T$, the number of steps is
$\lceil \text{tokens}/(BT)\rceil$ (\texttt{code/projects/a1/budget.py} does the arithmetic), and
the cosine schedule's end step must be set to that number for \emph{each} run
(\cref{sec:optim:sched}). Comparing batch sizes at a fixed number of steps instead would give larger
batches more data.

\begin{figure}[t]
\centering
\begin{tikzpicture}[
  ph/.style={draw, rounded corners=2pt, minimum height=0.8cm, align=center, font=\scriptsize, line width=0.6pt},
  arr/.style={-{Stealth[length=1.8mm]}, line width=0.6pt}]
\node[ph, fill=fdata!18, draw=keygreen, text width=2.3cm] (a) at (0,0) {debug config\\(tiny model,\\1--2 min)};
\node[ph, fill=fcompute!18, draw=sysblue, text width=2.6cm] (b) at (3.3,0) {coarse LR sweep\\$\times\sqrt{10}$ grid,\\short budget};
\node[ph, fill=fcompute!18, draw=sysblue, text width=2.6cm] (c) at (6.8,0) {fine LR sweep\\around best, full\\budget, $\ge$1 divergent};
\node[ph, fill=fmemory!20, draw=napkin, text width=2.6cm] (d) at (10.3,0) {batch sweep\\(re-tune LR per $B$,\\equal tokens)};
\node[ph, fill=fcomm!16, draw=labpurple, text width=2.6cm] (e) at (3.3,-1.8) {ablations: one\\change each, same\\budget \& seed};
\node[ph, fill=fcomm!16, draw=labpurple, text width=2.6cm] (f) at (6.8,-1.8) {OWT run with\\TinyStories recipe};
\node[ph, fill=warnred!12, draw=warnred, text width=2.6cm] (g) at (10.3,-1.8) {time-capped\\leaderboard run};
\draw[arr] (a)--(b); \draw[arr] (b)--(c); \draw[arr] (c)--(d);
\draw[arr] (c.south) -- ++(0,-0.35) -| (e.north); \draw[arr] (e)--(f); \draw[arr] (f)--(g);
\end{tikzpicture}
\caption{A sensible order for the experiments. Every arrow carries forward the best settings
found so far; every box writes a run directory with config, JSONL metrics (step and wall
clock) and a line in your experiment log.}
\label{fig:proj1:plan}
\end{figure}

\paragraph{Learning rate.} Sweep on a coarse logarithmic grid with short runs, then refine around
the best value with full-budget runs (\cref{fig:proj1:plan}). Include at least one learning rate
high enough to diverge. The handout asks you to relate the divergence point to the optimum, as
in \cref{fig:optim:lrsweep}. Report curves against both steps and wall-clock time.
\code{llmsys.explog.plot\_runs} draws either from JSONL logs.

\paragraph{Batch size.} Sweep from very small to the memory limit, re-tune the learning rate for
each batch size (at least a three-point search), compare at equal tokens, and discuss both loss
and wall-clock efficiency. \Cref{sec:optim:debug} gives the critical-batch-size intuition and
\cref{ch:scalinglaws} the scaling rules.

\paragraph{Ablations.} Change one thing at a time, keep the budget, seed and data order fixed, and
re-tune the learning rate when the change plausibly shifts it (removing normalisation does).
\Cref{ch:archdesign} explains what each ablation probes and Lab~7 provides a switchable
architecture on which to rehearse the methodology at toy scale. For the SwiGLU-vs-SiLU
comparison, match parameter counts as the handout specifies; otherwise you are measuring width,
not gating.

\paragraph{Logging and the report.} Every run gets a directory with its config, metadata (git
revision, hardware, versions) and metrics with step and wall time (\cref{sec:optim:explog}).
Keep a running experiment log: hypothesis, change, command, outcome, conclusion. Plots should label
the axes with units (nats/token), state the token budget, and use consistent colours for the same
setting across figures.

% =====================================================================
\section{OpenWebText and the leaderboard}\label{sec:proj1:owt}\covers{CS336 A1 \S7.4--7.5}

\paragraph{Interpreting the gap.} The same model and token budget give a much higher loss on
OpenWebText than on TinyStories. Loss in nats per token is not comparable across datasets (or across
tokenizers): it measures the model's cross-entropy relative to the data's inherent entropy, and web
text is far more diverse than synthetic children's stories. A useful framing is
$\text{loss} = \text{irreducible entropy} + \text{model and optimisation gap}$
(\cref{ch:scalinglaws}); OpenWebText raises both terms. Sampled text makes the difference
tangible. \Cref{ch:data} discusses what makes web data hard and how filtering changes it.

\paragraph{The leaderboard.}\label{sec:proj1:leaderboard} The task is to minimise validation
loss on OpenWebText within a fixed wall-clock budget on one GPU, using only the provided training
data. It is an optimisation problem over the whole system, and the book's chapters suggest the
axes to explore. We give directions, not a recipe:
\begin{itemize}
  \item \textbf{Throughput first.} Every percent of MFU is a percent more tokens. bf16 autocast,
        \code{torch.compile}, fused attention (SDPA/FlashAttention), a prefetching data loader, no
        per-step host syncs, and evaluation only at the end (\cref{ch:perf,sec:optim:data}).
  \item \textbf{Shape for the hardware.} Dimensions that are multiples of 64 or 128, a vocabulary
        padded to such a multiple, the largest batch that stays efficient, and a model size chosen
        for the budget rather than as large as possible (\cref{ch:scalinglaws,sec:archdesign:hparams}).
  \item \textbf{Optimisation.} Schedule shape (WSD or cosine ending exactly at the deadline),
        warmup, $\beta_2$, weight decay on the right parameter groups, and alternative optimizers
        such as Muon for hidden matrices (\cref{sec:optim:beyond}).
  \item \textbf{Architecture tweaks with good evidence at small scale}: weight tying (with an
        adjusted embedding initialisation), QK-norm, z-loss and other stability tricks that permit
        larger learning rates (\cref{sec:archdesign:stability}). The community ``speedrun''
        efforts for small GPT-style models document many such ideas; treat each one as a hypothesis
        to ablate.
  \item \textbf{Respect the clock.} Enforce the limit in your script (\texttt{wallclock.py} in
        \texttt{code/projects/a1/} stops training with time left for the final evaluation), and
        report the wall-clock axis the handout asks for.
\end{itemize}
Be aware that tricks that win a 45-minute contest do not necessarily help at larger scale; the
scaling-law project (\cref{ch:proj3}) is the antidote.

% =====================================================================
\section{Doing this outside Stanford}\label{sec:proj1:outside}

\begin{table}[h]
\centering\small
\begin{tabularx}{\linewidth}{@{}lX@{}}
\toprule
\textbf{Hardware} & \textbf{How to adapt}\\
\midrule
CPU only & Use the handout's low-resource token budget and target. Set \code{torch.set\_num\_threads}
           to your physical core count, use \code{torch.compile(model)}, and train BPE on the
           validation split first.\\
Apple silicon (MPS) & As above, with \code{device="mps"}; avoid TF32 on MPS and use the
           \code{aot\_eager} compile backend, as the handout recommends. Tokenizer work is CPU-bound
           either way.\\
Consumer GPU (12--24\,GB) & The full TinyStories configuration fits; expect several times the
           handout's B200 runtime. Use bf16 autocast if supported; reduce the batch size and
           accumulate gradients if needed (\cref{sec:pretrain:efficiency}).\\
Cloud H100/A100 & Close to the course setup. Rent by the hour, keep data on local NVMe,
           checkpoint to persistent storage, and script everything so a preempted instance can resume
           (\cref{sec:optim:ckpt}).\\
\bottomrule
\end{tabularx}
\caption{Hardware substitutions for Assignment~1.}
\end{table}

If you scale down, scale \emph{consistently}: fewer tokens and (if necessary) a smaller model, the
same experiment design, and the same logging. The qualitative findings (where the learning rate
diverges, which ablations hurt) usually survive the reduction; the absolute losses do not.

% =====================================================================
\section{Summary}
\begin{itemize}
  \item Assignment~1 builds a tokenizer, a modern Transformer LM, an optimizer and a training system
        from first principles, then studies them experimentally.
  \item The tests reach your code through thin adapters and compare numerics to snapshots, so
        follow the handout's conventions exactly (RoPE pairing, mask polarity, weight orientation).
  \item BPE at scale is a parallel map--reduce over chunks cut at special tokens, followed by a
        sequential, incrementally maintained merge loop. Streaming encoding must cut only where no
        token can straddle.
  \item Fix the token budget, derive steps from it, end the schedule there, and log step and
        wall-clock time for every run.
  \item For the leaderboard, measure and optimise throughput before you tune the model, and enforce
        the time limit in code.
\end{itemize}

\section{Exercises}

These exercises practise the guide's methods on examples that are \emph{not} assignment
problems.

\begin{exercise}[Pre-token boundaries]
Apply the GPT-2 pattern by hand to the string \verb*|x  = 10.5;  y| and list the pre-tokens.
Which alternative of the regex produced each one? Then check with the \code{regex} module.
\end{exercise}

\begin{exercise}[Chunk balance]
A 12\,GB file has documents whose lengths are roughly exponential with mean 4\,KB. You cut it
into $k$ chunks at document separators and process them with 8 workers. Estimate the expected
largest chunk-size deviation for $k=8$ and $k=64$, and explain why more chunks than workers help.
\end{exercise}

\begin{exercise}[Accounting method on another model]
Using \cref{sec:optim:account}, compute the parameter count and forward matmul FLOPs per token
for a pre-norm SwiGLU model with $V=32{,}000$, $T=2048$, $L=16$, $d=1024$, 16 heads and
$\dff=2816$, untied embeddings. What fraction of forward FLOPs is attention-score computation
($QK^\top$ and $PV$) at $T=2048$ and at $T=32{,}768$?
\end{exercise}

\begin{exercise}[Budgeted sweep design]
You have 6 GPU-hours, a 20-minute full-budget run, and 4-minute short runs that preserve the
learning-rate ranking. Design a learning-rate and batch-size study that respects the budget,
includes a divergent run, and re-tunes the learning rate per batch size.
\end{exercise}
